\fancyhead[L]{IHM21}
\renewcommand{\arraystretch}{1}


\titre{}{IHM21 - Interface client - serveur}


 \section{Définition}
 
 
 \cadreR{\defi  L'interaction entre un ordinateur \textbf{client} et un ordinateur \textbf{serveur} est à la base d'Internet. L'idée est de \textbf{permettre à une machine d'accéder à une ressource disponible sur une autre machine du réseau}. Une ressource peut aussi bien être une page web, qu'une base de données ou bien encore de la puissance de calcul. Il est important de comprendre ce qui est exécuté sur votre ordinateur et ce qui est exécuté sur le serveur. }
 
 \medskip
 
  Dans le cas extrême du « cloud-gaming » ou « cloud-working », toute la puissance de calcul est déportée vers le serveur. L'ordinateur client sert alors seulement d'interface. Il gère l'affichage et les périphériques : manette de jeux, souris, capteurs de réalité virtuelle. Il capte les événements et les transmet au serveur qui calcule et renvoie le résultat à afficher. 
  
  \medskip
  
 
\begin{liste}
	\item L'avantage est que l'ordinateur client (console de jeu, ...) n'a plus besoin d'une grande puissance de calcul.
	\item L'inconvénient est que cela nécessite une connexion internet de très bonne qualité. 
\end{liste}


\medskip

\begin{center}
\includegraphics[scale=0.45]{ScreenShot_20240704173424.ps}
\end{center}

Pour communiquer entre eux, le client et le serveur doivent parler la même langue c'est-à-dire utiliser les mêmes \textbf{protocoles}.

Pour les requêtes web, c'est le \textbf{protocole HTTP (HyperText Transfer Protocol)} qui est utilisé. 


\begin{center}
\includegraphics[scale=0.7]{client-serveur.ps}
\end{center}


\newpage
\section{Exécution côté client ou côté serveur?}

 Voici un exemple qui permet de comprendre la différence entre une action réalisée côté client et une autre côté serveur. Notre objectif est d'afficher l'heure. 
 
 \subsubsection{Première méthode: côté client en javascript}
 
 Voici le code d'une horloge réalisée en javascript
 
 
 \includegraphics[scale=0.5]{ScreenShot_20240704172816.ps}

 Cela fonctionne bien car tout s'exécute sur votre ordinateur (côté client), et il est inutile de recharger la page en envoyant une nouvelle requête au serveur. 
 
 \subsubsection{Deuxième méthode: côté serveur en PHP}
 
  Voici ce code tel qu'il est sur le serveur. On repère le code PHP car il est dans la balise : \verb+<?php ... ... ?>+.
  
  \includegraphics[scale=0.5]{ScreenShot_20240704172834.ps} 
  
On constate que l'horloge est bloquée. Afin d'obtenir une actualisation de l'heure, il faut réinterroger le serveur en réactualisant la page (touche F5).\\
Le serveur va réinterpréter le code PHP de la page et le renvoyer à vous : le client.\\

Il est possible de forcer l'actualisation de la page en utilisant la balise \verb+<meta>+ dans l'en-tête \verb+<head>+ :

\begin{lstlisting}[language = HTML]
  <meta http-equiv="Refresh" content="durée_en_secondes;URL=adresse_de_la_page">

\end{lstlisting}

C'est ce qui a été réalisé sur la page : \href{https://raisintine.fr/chocolatine/assets/P-08-client-serveur/horloge-php.php}{horloge-php.php}


\medskip

\rem Malgré la demande de rafraîchissement de la page, inclue dans le code de celle-ci, les méthodes d'optimisation des navigateurs pour diminuer la consommation de données et permettre un chargement plus rapide des pages peut rendre inefficace cette méthode.

Cela montre que pour ce genre d'application, il est préférable de déporter le calcul côté client (javascript) et ainsi soulager le réseau et le serveur. 

\medskip

\exer Depuis votre navigateur accéder aux Outils de développement (dans « Plus d'outils » pour Chrome et « Développement web » pour Firefox ou avec le raccourci CTRL+MAJ+I) et explorer le code de cette page à la recherche de code PHP


\section{Requêtes HTTP}

\subsubsection{Des requêtes différentes}

Entrer une adresse \textbf{URL (Uniform Ressource Locator)} dans la barre d'adresse de votre navigateur consiste à faire une \textbf{requête HTTP} à un \textbf{serveur distant} en lui demandant de fournir une ressource.
Voici quelques URL non-fonctionnelles juste pour l'exemple :

\begin{enumerate}
	\item \url{http://deodat.free.fr/}
    
    Par défaut, c'est la page index.html du site internet deodat hébergé sur le serveur free.fr qui sera envoyé par le serveur.
    \item \url{http://deodat.free.fr/NSI.php#test3}
    
    URL de la page \verb+NSI.php+ et plus précisément sur l'ancre \verb+test3+ de cette page. Une ancre correspond à la balise HTML ayant cette valeur d'id.
    
    \item \url{http://deodat.free.fr/NSI.php?nom=bob&test=3}
    
    URL de la page \verb+NSI.php+ avec les paramètres nom et test ayant pour valeurs bob et 3.
    
    Ces valeurs peuvent ensuite être récupérées et utilisées par le moteur PHP du serveur pour générer le code html de la page.
\end{enumerate}
    
    
    \medskip
    
    Le protocole d'échange est toujours indiqué au début de l'URL. Le protocole le plus utilisé est le HTTP.
    
Le \textbf{HTTPS} est un protocole d'échange \textbf{sécurisé}, c'est-à-dire que les données qui transitent entre le client et le serveur sont chiffrées.
 Veillez à vérifier que vous êtes en HTTPS lorsque vous effectuer des transactions bancaires en ligne.

\medskip

Voici un exemple de requête HTTP :

\includegraphics[scale=0.45]{ScreenShot_20240704172848.ps}

 Une requête HTTP contient des informations sur la \textbf{ressource demandée} mais aussi sur l'\textbf{utilisateur à l'origine de la requête}. 
 
 \medskip
 
\exer  
 
 \begin{enumerate}
 	\item     Depuis votre navigateur accéder aux Outils de développement (dans « Plus d'outils » pour Chrome et « Développement web » pour Firefox ou avec le raccourci CTRL+MAJ+I) ;
    \item Sélectionner l'onglet « Réseau » ou « Network » ;
    \item Observer les requêtes GET.


 \end{enumerate}
 
 
 
\import{../../../../preambule}{raz_compteurs.tex}
\newpage